% 本文件由 LaTeX 排版 Agent 根据有效 translation.json 自动生成。
% author=Victorinus; work_id=0711; title=论世界的创造

\chapter{论世界的创造}

% structure: p0001 source=h1 target=omitted reason=page-title-already-rendered-by-parent

当我默想并思虑我们被困于其中的这个世界的创造时，其创造之迅速，正如梅瑟的书（他关于其创造所写的，称之为\WorkTitle{创世纪}）所记载的。天主在六天内为祂的尊威产生了整个质料；第七天，祂以祝福祝圣了它。因着此缘故，由于在天上和地上的事物都按七日的数目安排，我将以涉及七的一切事物之原理之后，思考这第七日，并尽我所能，努力描绘这神圣能力的日子直到那圆满终结。\par

\BibleQuote{在第四日，祂在天上造了两个光体，较大的和较小的，使一个管昼，另一个管夜，}\BibleRef{创 1:16}——日月光体，又将其余星辰安置在天上，使它们照耀大地，并按它们的位置区分季节、年份、月份、日子和时辰。\par

如今显示了真理的理由：为何第四日被称为\OriginalTerm{第四日}{Tetras}，为何我们禁食直到第九时辰，甚至直到晚上，或者为何应逾越至次日。因此，我们这个宇宙由四元素构成——火、水、天、地。故此，这四元素构成了季节的\OriginalTerm{四重}{quaternion}。太阳和月亮也构成一年四季——春、夏、秋、冬；这些季节构成一个四重。再进一步从该原理出发，看，在天主宝座前有四活物，四福音，四条河流在乐园中流淌；\BibleRef{创 2:10}从亚当到诺厄的四代人，从诺厄到亚巴郎，从亚巴郎到梅瑟，从梅瑟到主基督、天主之子；还有四活物，即人、牛犊、狮子、鹰；四条河流：\Place{丕雄}、\Place{基红}、\Place{底格里斯}和\Place{幼发拉底}。那造生我们以上所述之事的本源——人基督耶稣，被恶人之手捉拿，由四名兵士看守。因此，因祂被四名兵士囚禁，因祂工作的庄严——为了使季节对人类健康、为丰收喜悦、为风暴平静得以运转——因此我们将第四日定为站礼或额外的斋戒。\par

第五日，陆地和水产了它们的后代。第六日，缺乏之物被创造；如此，天主从泥土中造了人，作为祂在陆地和水上创造的一切之主人。然而，祂在造人之前创造了天使和总领天使，将属灵的存在置于属土的存在之前。因为光在天空和大地之前被造。这第六日被称为\OriginalTerm{预备日}{parasceve}，即预备国度之日。因为祂成全了亚当，是按祂的肖像和模样造的。但祂之所以在创造天使和造人之前完成祂的工作，是免得他们错误地声称曾是祂的助手。这一日，因着主耶稣基督的苦难，我们要向天主做站礼或守斋。第七日，祂歇息了祂的一切工作，并祝福了它，祝圣了它。前一天，我们惯于严格守斋，以便在主的日子我们带着感恩前去领受我们的饼。并且，让预备日成为严格守斋的日子，免得我们似乎与犹太人一同遵守任何安息日，安息日的主基督自己曾藉先知说：\BibleQuote{祂的灵魂恼恨；}\BibleRef{依 1:13}这安息日祂亲自在身体上废弃了，然而，祂先前曾命令梅瑟，割礼不可逾越第八日，而这第八日常常落在安息日，正如我们在福音中所读到的。\BibleRef{若 7:22}梅瑟预见那人民的硬心，便在安息日举起双手，如此象征性地将自己钉在十字架上。\BibleRef{出 22:9}在战斗中，他们于安息日被外邦人追捕，以便被掳，并仿佛因法律的严格，被塑造以避免其教训。\BibleRef{加上 2:31}\par

因此，在第八日的第六篇圣咏中，\Person{达味}求主不要发怒责斥他，也不要在怒气中惩罚他；因为这正是那未来审判的第八日，它将超越七重秩序的界限。\Person{若苏厄}，\Person{纳维}的儿子，\Person{梅瑟}的继承人，他自己也打破了安息日；因为在安息日，他命令以色列子民\BibleRef{苏 6:4}拿着号角绕\Place{耶里哥}城墙行走，并向异族人宣战。\Person{玛塔提雅}，犹大的首领，也打破了安息日；因为在安息日，他杀死了\Person{叙利亚}王\Person{安提约古}的总督，并追击制服了外邦人。而且在\BibleBook{玛窦}{福音}中我们读到，\Person{依撒意亚}和他的同伴们也打破了安息日\BibleRef{玛 12:5}——为的是那真实而公义的安息日要在第七个千年时期中被遵守。因此，主赋予那七天各一千年；因为这样警告说：\Quote{主啊，在祢眼中，千年如同一日。}因此，在主的眼中，每个千年都是命定的，因为我发现主的眼目是七\BibleRef{匝 4:10}。因此，如我所叙述的，那真实的安息日将在第七个千年时期中，那时\Person{基督}将与祂的选民一同为王。此外，七重天也与那些日子相符；因为我们这样受到警告：\Quote{因主的话，诸天造成；因祂口中的神，万象造成。}有七种神。它们的名字就是那安息于\Person{天主}之\Person{基督}上的神，正如\Person{依撒意亚}先知所启示的：\Quote{智慧和聪敏之神，超见和刚毅之神，智慧和孝爱之神，以及敬畏上主之神，将安息在祂身上。}\BibleRef{依 11:2}因此，最高的天是智慧之天；第二天，聪敏之天；第三天，超见之天；第四天，刚毅之天；第五天，知识之天；第六天，孝爱之天；第七天，敬畏上主之天。因此，雷声轰鸣，闪电燃起，火焰积聚；火箭出现，星宿闪烁，由可怕彗星引起的焦虑被激起。有时太阳和月亮彼此靠近，引起那些不止是恐怖的现象，在其视域中放射光芒。但整个创造的作者是\Person{耶稣}。祂的名字是\OriginalTerm{圣言}{the Word}；因为祂的父这样说：\Quote{我的心发出美好的言语。}福音作者\Person{若望}这样说：\Quote{在起初已有圣言，圣言与\Person{天主}同在，圣言就是\Person{天主}。这圣言在起初就与\Person{天主}同在。万有是藉着祂而造成的；凡受造的，没有一样不是藉着祂而造成的。}因此，首先被造的是创造，其次是人类，人类的主宰——如宗徒所说。\BibleRef{格前 15:45}因此，这圣言，当它造光时，被称为智慧；当它造穹苍时，被称为聪敏；当它造陆地与海洋时，被称为超见；当它造日月及其他发光体时，被称为刚毅；当它召唤陆地和海洋时，被称为知识；当它造人时，被称为孝爱；当它祝福并圣化人时，便拥有敬畏上主之名。\par

看，羔羊的七角，\BibleRef{默 5:6} \Person{天主}的七眼\BibleRef{匝 4:10}——七眼就是羔羊的七神；\BibleRef{默 4:5} 在宝座前燃烧的七支火炬\BibleRef{默 4:5} 七盏金灯台，\BibleRef{默 1:13} 七只羔羊，\BibleRef{肋 23:18} \Person{依撒意亚}书中的七个女人，\BibleRef{依 4:1} \Person{保禄}书信中的七教会，七位执事，\BibleRef{宗 6:3} 七位天使，七支号角，书卷上的七个封印，完成五旬节的七个七日周期，\Person{达尼尔}书中的七周，以及\Person{达尼尔}书中的四十三周；与\Person{诺厄}一起，方舟中每种洁净的动物各七对；\Person{加音}的七倍报复，\BibleRef{创 4:15} 债务被豁免的七年，\BibleRef{申 15:1} 有七个口的灯，\BibleRef{匝 4:2} \Person{撒罗满}殿中的七根智慧之柱。\BibleRef{箴 11:1}\par

因此，现在你们可以看到，你们被告知关于\Person{天主}眷顾中无谬的光荣；但按我微小的能力所及，我将尽力阐述它。祂为藉着周再造那个\Person{亚当}，并给祂的全部受造界带来援助，这藉着祂的子，我们的主\Person{耶稣基督}的诞生而完成。那么，谁受了\Person{天主}法律的教导，谁充满了\Person{圣神}，心里不看到，在龙诱惑\Person{厄娃}的同一天，天使\Person{加俾额尔}向\Person{童贞玛利亚}报喜讯；在同一天，\Person{圣神}充满\Person{童贞玛利亚}，就是祂造光的那一天；在那一天，祂降生成人，就是祂造陆地和水的那一天；在同一天，祂被哺乳喂养，就是祂造星宿的那一天；在同一天，祂受割损，就是陆地和水产生它们产物的时候；在同一天，祂降生成人，就是祂用尘土造人的那一天；在同一天，\Person{基督}诞生，就是祂造人的那一天；在那一天，祂受难，就是\Person{亚当}堕落的那一天；在同一天，祂从死者中复活，就是祂造光的那一天？祂还在数字七中完成祂的人性：祂的诞生、婴儿期、童年期、少年期、青年期、壮年期、死亡。我也这样向犹太人陈述祂的人性：祂饥饿，祂口渴；祂赐予饮食；祂行走，又退隐；祂在枕头上睡眠；\BibleRef{谷 4:38} 此外，祂用脚行走在波涛汹涌的海上，祂命令风，祂医治病人，使瘸子复原，祂以言语使盲人看见——看，祂向他们宣告祂是主。\par

正如我上面所说的，白天被数字十二——白天和黑夜的十二个小时——分成两部分；而且，通过这些小时，月、年、季节和时代也被计算。因此，毫无疑问，也指定了十二个白天的天使和十二个黑夜的天使，即按照小时的数量。因为这些是坐在天主宝座前的白天和黑夜的二十四位见证人，他们头上戴着金冠，宗徒兼福音作者若望的\BibleBook{若望}{默示录}称他们为长老，原因是他们比其他天使和人都更年长。\par

\SourceRecord{Translated by Robert Ernest Wallis.；Ante-Nicene Fathers；Vol. 7.；Edited by Alexander Roberts, James Donaldson, and A. Cleveland Coxe.；Buffalo, NY: Christian Literature Publishing Co.,；1886.；<http://www.newadvent.org/fathers/0711.htm>.}
